% 本文件由 tools/build_book.py 从 dist/reports/assertions.json 生成，请勿手改。
\begin{longtable}{@{}p{0.06\textwidth}p{0.32\textwidth}p{0.22\textwidth}p{0.16\textwidth}@{}}
\caption{章节 4.5 数值断言（生成时间 2026-10-18T17:31:25-0400）}\\
\toprule
类别 & 断言 & 理论值 & 实测值 \\
\midrule
\endfirsthead
\toprule
类别 & 断言 & 理论值 & 实测值 \\
\midrule
\endhead
\bottomrule
\endlastfoot

B & 帧角点：有限性 + requiredAux + 掩码自洽 & 数组与 meta 全有限、\allowbreak{}能编 .nlf；requiredAux 命中；地板/\allowbreak{}占位/\allowbreak{}哨兵掩码与数据一致 & \{'first': [],\allowbreak{} 'n\_frames': 104,\allowbreak{} 'n\_problems': 0\} \checkmark \\
B & meta 与数组的自洽性(四页) & meta 的每个结论都能从同帧数组复算(\allowbreak{}平滑因子/\allowbreak{}地板计数/\allowbreak{}因子统计/\allowbreak{}损失与别名) & \{'n\_problems': 0,\allowbreak{} 'problems': []\} \checkmark \\
B & 纯函数与时间钳制 & 同参数两次构建逐位相同；scrub 在 [0,\allowbreak{}tMax] 外钳制；mode=\allowbreak{}none 不声明 tMax/\allowbreak{}dt/\allowbreak{}fps & \{'n\_problems': 0,\allowbreak{} 'problems': []\} \checkmark \\
C & 跨页阈值：高频消失、低频不动、MG 与平滑子的量级差 & 10 步后高频 <=\allowbreak{}1e-4 而低频 >=\allowbreak{}0.99；MG 总代价比平滑子 >=\allowbreak{}100 倍；纯高频残差的不可见比例 >=\allowbreak{}0.9 且混合残差损失至少掉 10 倍 & \{'viz1\_high\_survival': 1.686706566226433e-05,\allowbreak{} 'viz1\_low\_survival': 0.9979938376613837,\allowbreak{} 'viz3\_cost\_ratio': 800.3068181818182,\allowbreak{} 'viz3\_mgcg\_iters': 7.0,\allowbreak{} 'viz3\_two\_grid': 0.11111111111111113,\allowbreak{} 'viz4\_mixed\_drop': 390.82270979700394,\allowbreak{} 'viz4\_pure\_loss': 0.9999851987624295,\allowbreak{} 'viz4\_pure\_norm': 0.0012029456640646609\} \checkmark \\
B & manifest 自洽性 & 参数边界/\allowbreak{}默认值/\allowbreak{}临界值/\allowbreak{}枚举自洽，轴 >=\allowbreak{} 2 且 autoscaled 显式，pitfalls 与 theory 非空 & \{'n\_problems': 0,\allowbreak{} 'problems': []\} \checkmark \\
A & 转移恒等式最大偏差 & max|R-0.5P\textasciicircum{}T|=\allowbreak{}0，max|R A\_h P-A\_H|=\allowbreak{}0，max|P\textasciicircum{}T A\_h P-2A\_H|=\allowbreak{}0 & \{'PtAP\_minus\_2AH': 0.0,\allowbreak{} 'R\_minus\_half\_PT': 0.0,\allowbreak{} 'galerkin\_defect': 0.0,\allowbreak{} 'poisson\_vs\_1\_5': 0.0\} \checkmark \\
A & Jacobi 逐模态因子：实测 vs 理论 vs 1.5 章口径 & |1-omega(\allowbreak{}1-cos theta\_k)|；omega=\allowbreak{}2/\allowbreak{}3 时 |1-(\allowbreak{}4/\allowbreak{}3)sin\textasciicircum{}2(\allowbreak{}theta\_k/\allowbreak{}2)| & \{'measured\_vs\_theory': 4.447949739383039e-15,\allowbreak{} 'theory\_vs\_closed\_form': 2.220446049250313e-16\} \checkmark \\
A & smoother\_\allowbreak{}matrix vs 1.5 章迭代矩阵 & 与 jacobi/\allowbreak{}gs/\allowbreak{}sor\_iteration\_matrix 逐位相同(\allowbreak{}omega 由 I-omega(\allowbreak{}I-T\_J) 缩放) & 0 \checkmark \\
A & 两网格因子 vs 1/9 & rho\_TG =\allowbreak{} mu\_smooth\textasciicircum{}2 =\allowbreak{} (\allowbreak{}1/\allowbreak{}3)\textasciicircum{}2 =\allowbreak{} 0.111111... & \{'factors': [0.11111111111111104,\allowbreak{} 0.11111111111111067,\allowbreak{} 0.11111111111111116,\allowbreak{} 0.11111111111111077],\allowbreak{} 'max\_abs\_dev': 4.3021142204224816e-16,\allowbreak{} 'spread': 4.85722573273506e-16\} \checkmark \\
A & FAS vs 离散牛顿参考解 & FAS 不动点 =\allowbreak{} 离散方程的解：max|u\_FAS - u\_Newt| 到机器精度；残差 <=\allowbreak{} 1e-8 & \{'discrete\_vs\_exact': [0.0025975536252749176,\allowbreak{} 0.0006490957336637848,\allowbreak{} 0.00016225547489456815],\allowbreak{} 'exact\_error': [7.593925488436071e-14,\allowbreak{} 1.156852391659413e-13,\allowbreak{} 1.3233858453531866e-13],\allowbreak{} 'final\_residual': [1.3547766369373269e-11,\allowbreak{} 7.5482179359868e-11,\allowbreak{} 3.071030313191496e-10]\} \checkmark \\
B & 粗层校正对 P 值域的零化(相对) & R A\_h P =\allowbreak{} A\_H $\Rightarrow$ (\allowbreak{}I-P A\_H\textasciicircum{}\{-1\} R A\_h) P y =\allowbreak{} 0 & $3.082\times 10^{-14}$ \checkmark \\
B & Galerkin 投影 vs 两个假恒等式 & Q\textasciicircum{}2=\allowbreak{}Q、\allowbreak{}A-对称(\allowbreak{}偏差 0)；R P =\allowbreak{} 0.75I+0.125 adj(\allowbreak{}偏差 0)；max|RP-I| =\allowbreak{} 0.25 !=\allowbreak{} 0 & \{'projection\_defect': 4.4408920985006214e-15,\allowbreak{} 'rp\_closed\_form': 0.0,\allowbreak{} 'rp\_vs\_identity': 0.25\} \checkmark \\
B & 含逆交换律的反例留档 & 一般右端 rel$\approx$1.0；子空间上 abs$\approx$0.61；真命题 A\_H\textasciicircum{}\{-1\} R A\_h P =\allowbreak{} I 偏差 0 & \{'general\_rhs\_rel': 1.0098227246529399,\allowbreak{} 'subspace\_abs': 0.6126289540840444,\allowbreak{} 'true\_form\_abs': 1.7763568394002505e-15\} \checkmark \\
B & 循环算子线性性与矩阵一致性(相对) & E 线性：E(\allowbreak{}x1+x2,\allowbreak{}b1+b2)=\allowbreak{}E(\allowbreak{}x1,\allowbreak{}b1)+E(\allowbreak{}x2,\allowbreak{}b2)，且 E 的列就是 mg\_apply(\allowbreak{}e\_j) & $7.406\times 10^{-15}$ \checkmark \\
B & 计数器 vs 闭式代价(相对) & V: sum\_l 2*n\_l/\allowbreak{}n\_0 + 24*n\_L/\allowbreak{}n\_0；W: 重数 2\textasciicircum{}l & \{'rows': [[15,\allowbreak{} 'v',\allowbreak{} 4.933333333333334,\allowbreak{} 4.933333333333334],\allowbreak{} [15,\allowbreak{} 'w',\allowbreak{} 18.26666666666667,\allowbreak{} 18.266666666666666],\allowbreak{} [31,\allowbreak{} 'v',\allowbreak{} 4.387096774193548,\allowbreak{} 4.387096774193548],\allowbreak{} [31,\allowbreak{} 'w',\allowbreak{} 19.677419354838705,\allowbreak{} 19.67741935483871]],\allowbreak{} 'worst\_rel': 4.988580128479158e-16\} \checkmark \\
B & 估计器：对法 vs 错法 vs 显式矩阵谱半径 & 归一化后每步范数即放大率(\allowbreak{}几何平均 -> rho(\allowbreak{}E))；相邻相除会得到 $\approx$1 的假值 & \{'cycle\_factor\_api': 0.1837371257461596,\allowbreak{} 'cycle\_factor\_correct': 0.1837371257461596,\allowbreak{} 'naive\_ratio\_of\_ratios': 1.0037651335386966,\allowbreak{} 'rho\_of\_matrix': 0.18955483498701953\} \checkmark \\
B & injection 的伴随缺口与两网格因子代价 & R =\allowbreak{} 0.5 P\textasciicircum{}T(\allowbreak{}满权)才使 R A\_h P =\allowbreak{} A\_H；injection 下两网格因子显著变大(\allowbreak{}实测约 3 倍) & \{'fw': [0.11111111111111067,\allowbreak{} 0.11111111111111116,\allowbreak{} 0.11111111111111077],\allowbreak{} 'inj': [0.32906339564483605,\allowbreak{} 0.3322632725939292,\allowbreak{} 0.33306559645892087],\allowbreak{} 'min\_ratio': 2.961570560803536,\allowbreak{} 'worst\_asymmetry': 0.33333333333333337\} \checkmark \\
C & GS 高频放大率与边界误差收缩 & max\_high |g(\allowbreak{}theta)| -> 1/\allowbreak{}sqrt(\allowbreak{}5)=\allowbreak{}0.447214；边界误差 \textasciitilde{} O(\allowbreak{}h)(\allowbreak{}比值 -> 2) & \{'gap': [0.024552769515947204,\allowbreak{} 0.01224299769295889,\allowbreak{} 0.006093617540537721,\allowbreak{} 0.0030413549001271356],\allowbreak{} 'gap\_ratios': [2.0054540670270504,\allowbreak{} 2.0091509865055506,\allowbreak{} 2.0035864740030815],\allowbreak{} 'max\_abs\_dev': 0.0018672883037116894,\allowbreak{} 'max\_high': [0.43205008623077534,\allowbreak{} 0.4396968652810592,\allowbreak{} 0.443471156521669,\allowbreak{} 0.44534630719624624],\allowbreak{} 'young\_eigen': [0.9619397662556434,\allowbreak{} 0.9903926402016151,\allowbreak{} 0.9975923633360985,\allowbreak{} 0.9993977281025862]\} \checkmark \\
C & 缩放常数 c 的因子扫描 & |c-0.5|=\allowbreak{}0 时循环收敛；c>0.5 过校正可发散(\allowbreak{}粗层校正被放大 2c 倍) & \{'c': [0.25,\allowbreak{} 0.5,\allowbreak{} 0.75,\allowbreak{} 1.0,\allowbreak{} 2.0],\allowbreak{} 'cycle\_factor': [0.9345419144624127,\allowbreak{} 0.188663987137861,\allowbreak{} 4.343326780434276,\allowbreak{} 19.196681369289028,\allowbreak{} 549.1505837323252],\allowbreak{} 'galerkin\_defect': [0.5,\allowbreak{} 0.0,\allowbreak{} 0.5,\allowbreak{} 1.0,\allowbreak{} 3.0],\allowbreak{} 'worst': 549.1505837323252\} \checkmark \\
C & V vs W：因子、每周期代价、总代价 & W 因子接近两网格(\allowbreak{}更小)，但每周期代价约 2\textasciicircum{}L 级；总代价由 因子 与 代价 共同决定 & \{'cost\_ratio': 4.203331020124932,\allowbreak{} 'cost\_v': 45.74603174603173,\allowbreak{} 'cost\_w': 192.2857142857151,\allowbreak{} 'v\_cycles': 11,\allowbreak{} 'v\_factor': 0.188663987137861,\allowbreak{} 'v\_work': 4.158730158730157,\allowbreak{} 'w\_cycles': 9,\allowbreak{} 'w\_factor': 0.11448953348269528,\allowbreak{} 'w\_work': 21.365079365079456,\allowbreak{} 'work\_ratio': 5.137404580152696\} \checkmark \\
C & 因子扫描：MG 的平坦 vs 平滑子的 h\textasciicircum{}2 恶化 & MG 因子与 h 无关(\allowbreak{}std 小、\allowbreak{}漂移小)；1-rho\_smoother \textasciitilde{} h\textasciicircum{}2 $\Rightarrow$ 斜率 -2(\allowbreak{}渐近) & \{'drift': 0.02922987995958487,\allowbreak{} 'factor\_mg': [0.16093827456842116,\allowbreak{} 0.1837371257461596,\allowbreak{} 0.18935184490011966,\allowbreak{} 0.188663987137861,\allowbreak{} 0.19016815452800603],\allowbreak{} 'factor\_smoother': [0.9492530216741912,\allowbreak{} 0.9871901869343185,\allowbreak{} 0.9967898177764185,\allowbreak{} 0.9991909271994155,\allowbreak{} 0.9997104469939833],\allowbreak{} 'factor\_two\_grid': [0.11111111111111113,\allowbreak{} 0.11111111111111104,\allowbreak{} 0.11111111111111067,\allowbreak{} 0.11111111111111116,\allowbreak{} 0.11111111111111077],\allowbreak{} 'iters\_mg': [2.0,\allowbreak{} 2.0,\allowbreak{} 2.0,\allowbreak{} 2.0,\allowbreak{} 2.0],\allowbreak{} 'iters\_smoother': [56.0,\allowbreak{} 226.0,\allowbreak{} 907.0,\allowbreak{} 3600.0,\allowbreak{} 10061.0],\allowbreak{} 'mean': 0.18257187737611352,\allowbreak{} 'n': [7,\allowbreak{} 15,\allowbreak{} 31,\allowbreak{} 63,\allowbreak{} 127],\allowbreak{} 'smoother\_slope': -1.8891537362805895,\allowbreak{} 'std': 0.011046771127791241\} \checkmark \\
C & CG vs MGCG 迭代数(tol=1e-8) & CG \textasciitilde{} O(\allowbreak{}h\textasciicircum{}-1)(\allowbreak{}sqrt(\allowbreak{}kappa))；MGCG 迭代数 O(\allowbreak{}1)；加速比按迭代数 \textasciitilde{} 1/\allowbreak{}h & \{'cg\_iters': [15,\allowbreak{} 31,\allowbreak{} 63,\allowbreak{} 127,\allowbreak{} 255],\allowbreak{} 'mgcg\_iters': [7,\allowbreak{} 7,\allowbreak{} 7,\allowbreak{} 7,\allowbreak{} 7],\allowbreak{} 'mgcg\_std': 0.0,\allowbreak{} 'n': [15,\allowbreak{} 31,\allowbreak{} 63,\allowbreak{} 127,\allowbreak{} 255],\allowbreak{} 'speedup': [2.142857142857143,\allowbreak{} 4.428571428571429,\allowbreak{} 9.0,\allowbreak{} 18.142857142857142,\allowbreak{} 36.42857142857143]\} \checkmark \\
C & FAS 因子(跨 h)与线性化反例 & FAS 因子与 h 无关；线性化残差收敛到 1e-10 时真残差仍 O(\allowbreak{}1)，假解误差 \textasciitilde{}8e-2(\allowbreak{}真解幅度 1) & \{'false\_linear\_residual': 4.075381651270153e-10,\allowbreak{} 'false\_solution\_error': 0.07886428380016719,\allowbreak{} 'false\_true\_residual': 5.5964906255189675,\allowbreak{} 'fas\_factor': [0.15442316308010817,\allowbreak{} 0.1734676284798888,\allowbreak{} 0.18617952604063434,\allowbreak{} 0.18917931305791433],\allowbreak{} 'gap': 13732433191.31777,\allowbreak{} 'mean': 0.17581240766463643,\allowbreak{} 'n': [15,\allowbreak{} 31,\allowbreak{} 63,\allowbreak{} 127],\allowbreak{} 'std': 0.013685293900176134\} \checkmark \\
B & 窗口因子 vs 触底后的朴素比值 & 触底段比值 \textasciitilde{}1；窗口估计器跳过地板与初始瞬态后应给出真衰减率 0.19 & \{'naive\_tail\_ratio': 1.0,\allowbreak{} 'window': [7,\allowbreak{} 11],\allowbreak{} 'window\_factor': 0.19\} \checkmark \\
\end{longtable}
